\section{Introduction to Merging}\label{section:merging_introduction}

In Bayesian statistics, learning about the data-generating mechanism begins with the specification of a prior distribution over a class of probability distributions that could have generated the data. This prior may be uninformative (e.g., uniform over a suitable class), reflecting limited prior information, or more concentrated when substantive prior knowledge is available. Upon observing data, the statistician updates these beliefs via Bayes' rule, obtaining a posterior distribution over the candidate data-generating laws.

The goal of this section is to compare posterior beliefs arising from different prior distributions. This problem is commonly referred to as the merging of opinions. To study it rigorously, several modeling choices must be clarified. In particular, what assumptions are made about the sequence of observations? Are restrictions imposed on the class of prior distributions?

% In finite-dimensional parametric settings, aspects of this problem have been investigated; see \cite{ley2017distances}. The authors use a distance-based approach and conclude that merging occurs when prior distributions are defined over finite-dimensional parameters.

This problem was studied in \cite{blackwell1962merging}, where the convergence of predictive distributions is analyzed. Their framework considers probability measures defined on the space of infinite sequences of observations. After observing the first $n$ observations, one forms a predictive distribution over future observations via conditional probabilities. They show that if one probability measure $Q$ is absolutely continuous with respect to another probability measure $P$, then their corresponding predictive distributions merge almost surely with respect to $Q$. This result applies to Bayesian settings in which priors are defined over infinite-dimensional spaces and implies that the corresponding predictive distributions, and hence posterior beliefs, merge. However, it relies on a key restriction: one probability distribution on the sequence of infinite observations induced by a prior must be absolutely continuous with respect to the other. A sufficient condition guaranteeing this is to impose absolute continuity at the level of prior distributions. However, this condition can fail even in standard Bayesian nonparametric models; for example, two Dirichlet process priors with mutually singular base measures may themselves be mutually singular. Moreover, the result is formulated under the assumption that the observations are generated according to one of the probability measures under consideration. If the objective is instead to compare posterior updating procedures themselves, it is natural to avoid strong assumptions about the true data-generating mechanism.

This perspective is adopted in \cite{catalano2026merging}, and the goal of this chapter is to extend the study of merging to another example of a Bayesian nonparametric model.
% \textcolor{red}{Missing: consequence of no assumption on the data-generating mechanism. Merging might occur, but not consistency? Merging might not occur at all if the data are not generated from the prior-likelihood law.}

